\begin{textsample}{POS Dim 3 – generic\_gpt – Score 2.00 – generic\_gpt/t036\_gpt.txt}  \label{ex:f3_pos_036}
I don't know how to explain this without sounding dramatic, but lately it feels like I ’m disappearing from my own life.
I ’m not completely isolated in the obvious sense.
I have coworkers I chat with, a family group chat that occasionally pings, a couple of old \textbf{friends} I exchange memes with.
On paper, it looks like I ’m “ connected. ” But none of it feels like it actually reaches me.
It ’s like I ’m standing behind glass, watching myself go through the motions of being a person among other people.
The worst part is how ordinary everything looks from the outside.
I go to work, I answer emails, I say “ good weekend? ” on Mondays and “ have a good one ” on Fridays.
I laugh at the right \textbf{times}.
But underneath, there ’s this constant ache that nobody really knows me, and I ’m not sure I even remember how to let them.
I used to have deeper friendships, but they faded slowly—different cities, different schedules, people pairing off and having kids.
I can't point to a single moment where I “ lost ” anyone ; it was more like we all just stopped being central to each other ’s lives.
Now, every \textbf{time} I think about reaching out, I feel weirdly ashamed, like I ’m trying to resurrect something that quietly died \textbf{years} ago.
I try to be proactive—join \textbf{things}, say yes to invitations when they appear—but conversations never seem to get past the surface.
I leave social \textbf{situations} feeling more alone than before, because it reminds me how far I am from the \textbf{kind} of closeness I actually want.
Most days I just come home, sit in the quiet, scroll through other people ’s lives, and wonder when exactly mine started feeling like a waiting room.

% matched lemmas: friend, kind, situation, thing, time, year
\end{textsample}
